% \appendix
\setcounter{secnumdepth}{1}

% \renewcommand{\thesection}{Appendix \Alph{section}}
% \setcounter{section}{0}
% \setcounter{figure}{0}
% \setcounter{table}{0}
% Redefine numbering format
% \renewcommand{\thefigure}{S\arabic{figure}}
% \renewcommand{\thetable}{S\arabic{table}}
\renewcommand{\thetable}{A\arabic{table}}
\setcounter{table}{0}

% \onecolumn
\begin{center}
    {\LARGE \bfseries Supplementary Material}
\end{center}
\vspace{1em} 

This supplementary material provides detailed information organized as follows:
\begin{itemize}
    \item \textbf{Appendix~\ref{apd:full_tables}}: Complete forecasting performance tables for all models on DeepSeek-R1 and Llama-3.3-70B annotations
    \item \textbf{Appendix~\ref{apd:masking_l33_dsr1}}: Sensitivity analysis of masking history on Llama-3.3-70B annotations
    \item \textbf{Appendix~\ref{apd:encoder-forecasting}}: Technical details of encoder-based forecasting methods, including MLP-head and masking approaches
    \item \textbf{Appendix~\ref{apd:encoder-survival}}: Survival modeling framework using language model embeddings and the hyperparameters used
    \item \textbf{Appendix~\ref{apd:compute_resource}}: Computing infrastructure used for experiments, including hardware specifications, memory requirements, operating system, and software environment
    \item \textbf{Appendix~\ref{TTS_prompt}}: LLM prompt template for extracting clinical time series from case reports
    \item \textbf{Appendix~\ref{apd:prompts}}: Prompting strategies for zero-shot and few-shot forecasting
    \item \textbf{Appendix~\ref{apd:open_source}}: Analysis of open-source models to address data leakage concerns
    \item \textbf{Appendix~\ref{apd:broader_impact}}: Broader societal impact and ethical considerations
\end{itemize}

\section{Forecasting performance on \textit{DeepSeek-R1} and \textit{Llama-3.3-70B} textual time-series annotations}
\label{apd:full_tables}

\input{AAAI_2026/Tables/DS_time_ordered}
\input{AAAI_2026/Tables/L33_time_ordered}

Tables~\ref{tab:DS_forecasting} and~\ref{tab:L33_forecasting} present the forecasting performance of our models on textual time-series annotations generated by \textit{DeepSeek-R1} and \textit{Llama-3.3-70B}, respectively. We report standard forecasting metrics across multiple evaluation settings to provide a comprehensive view of model performance. 

%%%%%%%%%%%%%%%%%%%%%%%%%%%%%%%%%%%%%%%%%%%%%%%%%%%%%%%%%%%%%%%%%
%%%%%%%%%%%%%%%%%%%%%%%%%%%%%%%%%%%%%%%%%%%%%%%%%%%%%%%%%%%%%%%%
%\clearpage
\section{Masking history: \textit{Llama-3.3-70B} textual time-series annotations}
\label{apd:masking_l33_dsr1}

Table~\ref{tab:masking_effects_llama} examines the impact of randomly masking portions of the input history on model performance when using \textit{Llama-3.3-70B}-generated annotations. 
We evaluate performance across three datasets (T2S2-L33, Sepsis-10, and Sepsis-100) using F1 scores for prediction of next \(k\) events happening with the next day and concordance index (c-index) correctly ordering the next \(k\) events. 
As the proportion of masked history increases, we observe a general decline in F1 performance, particularly at higher masking rates (e.g., 90\%), suggesting that access to historical context is important for accurate forecasting. 
Interestingly, the concordance index remains relatively stable, indicating the model's robustness in ranking the next \(k\) events with partial input histories.

\begin{table*}[!htbp]
    \centering
    \caption{Effect of Randomly Masking History (\textit{Llama-3.3-70B} annotations) on F1 and Concordance. \textbf{Abbreviations}: F1(1d)-10/100: F1 (1 day) for sepsis-10 and sepsis-100 respectively. c10/100: c-index for sepsis-10 and sepsis-100 respectively.}
    \resizebox{0.65\linewidth}{!}{
    \begin{tabular}{lcccccc}
        \toprule
        \textbf{Timestep drop rate} & F1(1d) & c-index & F1(1d)-10 & c10 & F1(1d)-100 & c100 \\
        \hline
        0\%  & 0.645 & 0.632 & 0.431 & 0.595 & 0.626 & 0.612 \\
        15\% & 0.628 & 0.635 & 0.422 & 0.592 & 0.624 & 0.610 \\
        30\% & 0.634 & 0.634 & 0.423 & 0.586 & 0.621 & 0.598 \\
        60\% & 0.604 & 0.634 & 0.391 & 0.577 & 0.605 & 0.613 \\
        90\% & 0.531 & 0.622 & 0.328 & 0.546 & 0.534 & 0.607 \\
        \bottomrule
    \end{tabular}
    }
    \label{tab:masking_effects_llama}
\end{table*}

%%%%%%%%%%%%%%%%%%%%%%%%%%%%%%%%%%%%%%%%%%%%%%%%%%%%%%%%%%%%%%%%%
%%%%%%%%%%%%%%%%%%%%%%%%%%%%%%%%%%%%%%%%%%%%%%%%%%%%%%%%%%%%%%%%
%\clearpage
\section{Textual Time Series Forecasting Framework of Encoder Models}
\label{apd:encoder-forecasting}

We develop a general framework for predictive modeling over textual time series of the encoder-only models, designed to accommodate the irregular, event-based structure common to clinical case reports. 
As highlighted Section~\ref{sec:tts-annotation}, each data instance consists of a sequence of clinical events, where each event is a free-text string describing a clinical finding, intervention, or diagnosis, paired with a timestamp in hours relative to the time of admission. 
The format resembles a two-column table, where the first column contains textual clinical events and the second column contains integer-valued timestamps.

\paragraph{Input Preprocessing and Timeseries Construction.}
From each case report, we construct multiple training examples for sequence forecasting. 
For a fixed forecast horizon parameter $K$, we generate overlapping windows of historical events and forecast targets. 
Specifically, for each case report with $L$ rows, we construct a sequence of up to $L-K$ forecasting examples. 
Each example is structured as follows:

\begin{itemize}
    \item A \textbf{history segment} $\mathcal{H}_t = \{(e_i, t_i)\}_{i=1}^{n}$, consisting of clinical events $e_i$ occurring at timestamps $t_i \leq t$.
    \item A \textbf{forecast target} $\mathcal{F}_t = \{(f_j, \tau_j)\}_{j=1}^{K}$, consisting of the next $K$ events after $t$.
\end{itemize}

To preserve the temporal coherence of events, we treat all events with the same timestamp as an atomic unit. 
That is, if multiple events share the same timestamp, they are included together in the history or forecast block, and the next forecasting window is only advanced after the full timestamp block has been consumed.

\paragraph{Model Input Representation.}
Each training example is serialized into a transformer-compatible input by linearizing the historical events and forecast events. 
Importantly, we preserve both temporal and textual content by encoding each event in the form ``\texttt{t:e}'' where \texttt{t} is the timestamp and \texttt{e} is the event string.
The input to the language model takes the form:

\begin{equation}
\texttt{[CLS]} \; t_1\texttt{:}e_1 \; \texttt{[SEP]} \; t_2\texttt{:}e_2 \; \texttt{[SEP]} \; \dots \; \texttt{[SEP]} \; t_n\texttt{:}e_n
\end{equation}

This is followed by additional target-specific components, as described below.

\paragraph{Language Models.}
To test the generality of our framework, we experiment with a variety of pre-trained encoder-only language models, including:
\begin{itemize}
    \item \textit{bert-base-uncased} (BERT; Devlin et al., 2018)
    \item \textit{roberta-base} (Liu et al., 2019)
    \item \textit{microsoft/deberta-v3-small} (He et al., 2021)
    \item \textit{answerdotai/ModernBERT-base} and \textit{answerdotai/ModernBERT-large}
    \item \textit{thomas-sounack/BioClinical-ModernBERT-base} and \textit{thomas-sounack/BioClinical-ModernBERT-large} 
\end{itemize}
Each model is used with its corresponding tokenizer. 
For models with SentencePiece tokenization (e.g., DeBERTa), we use the appropriate fast tokenizer such as \textit{DebertaV2Tokenizer} to ensure compatibility. 
All inputs are truncated from the left (i.e., keeping the most recent tokens) if they exceed the model's context window (typically 512 tokens).

\paragraph{Prediction Tasks.}
We define two primary prediction tasks over $\mathcal{F}_t$, each formulated as a binary classification problem.

\subparagraph{1. Event Ordering (Concordance) Task.}
The goal in this task is to predict the correct temporal ordering between future events. 
For each example, we generate all \emph{comparable pairs} of events in $\mathcal{F}_t$, defined as pairs $(f_i, f_j)$ where $\tau_i \neq \tau_j$. 
For each pair, we form a binary classification example as follows:

\begin{itemize}
    \item \textbf{Input:} The serialized history segment followed by:
    \begin{equation}
        \texttt{[SEP]} \; \tau_i\texttt{:}f_i \; \texttt{[SEP]} \; \tau_j\texttt{:}f_j \; \texttt{[SEP]}
    \end{equation}
    \item \textbf{Label:} $y = 1$ if $\tau_i < \tau_j$, else $y = 0$
\end{itemize}

The model output is taken from the \texttt{[CLS]} token, passed through a linear classification head, and softmaxed to produce a binary probability. 
We optimize cross-entropy loss across all pairwise comparisons. 
At evaluation time, we report the concordance index (c-index), which measures the proportion of correctly predicted pairs among all comparable pairs.

\subparagraph{2. Time-Window Classification Task.}
This task evaluates whether each forecast event occurs within a fixed prediction window $H$ hours of the last historical event. 
For each $f_j \in \mathcal{F}_t$, we generate an input:

\begin{itemize}
    \item \textbf{Input:} The serialized history segment followed by:
    \begin{equation}
        \texttt{[SEP]} \; \tau_j\texttt{:}f_j \; \texttt{[SEP]}
    \end{equation}
    \item \textbf{Label:} $y = 1$ if $\tau_j - t_n \leq H$, where $t_n$ is the timestamp of the last historical event.
\end{itemize}

This is again treated as a binary classification task over the \texttt{[CLS]} embedding, using a linear layer and softmax. 
We report macro-averaged F1 scores over all prediction examples.

\paragraph{Masked Language Modeling Approaches.}
In addition to the MLP-head approach described above, we implement masked language modeling (MLM) variants that leverage the pretrained masked token prediction capabilities of encoder models. These approaches reformulate our forecasting tasks as cloze-style problems:

\subparagraph{1. Masked Event Occurrence Prediction.}
For the time-window classification task, we construct prompts that include a \texttt{[MASK]} token whose prediction indicates whether an event occurs within the specified horizon:

\begin{itemize}
    \item \textbf{Input:} History formatted as: \texttt{[CLS]} $h_1$ \texttt{[SEP]} $h_2$ \texttt{[SEP]} ... \texttt{[SEP]} Will "$e_j$" happen within $H$ hours? \texttt{[MASK]} \texttt{[SEP]}
    \item \textbf{Target:} The model predicts either "\texttt{yes}" or "\texttt{no}" at the \texttt{[MASK]} position
\end{itemize}

We constrain the output vocabulary to only these two tokens and compute cross-entropy loss over their logits. For fine-tuned variants, we train the entire model to optimize this masked prediction task. For zero-shot variants, we directly use the pretrained model's predictions without additional training.

\subparagraph{2. Masked Temporal Ordering.}
For the concordance task, we predict relational tokens that indicate temporal ordering:

\begin{itemize}
    \item \textbf{Input:} History followed by: \texttt{[SEP]} $e_1$ \texttt{[MASK]} $e_2$ \texttt{[SEP]}
    \item \textbf{Target:} The model predicts "\texttt{before}" if $e_1$ occurs before $e_2$, else "\texttt{after}"
\end{itemize}

Again, we restrict predictions to these two tokens. This formulation allows the model to leverage its pretrained understanding of temporal relationships encoded during MLM pretraining. 

Note that in our dataset, we randomize the relative ordering of $e_1$ and $e_2$ each time to avoid introducing positional biases that could lead the model to rely on superficial cues rather than genuine temporal reasoning.

\subparagraph{Implementation Details.}
For both tasks, we extract the logits at the \texttt{[MASK]} position and apply a linear projection to the target vocabulary size (2 tokens). The key advantage of this approach is that it can leverage pretrained masked language modeling capabilities, potentially requiring less task-specific training. We evaluate both fine-tuned versions (which update all model parameters) and zero-shot versions (which use frozen pretrained weights).

\paragraph{Training and Evaluation.}
Each task is trained separately using AdamW with linear learning rate decay and early stopping based on validation loss. 
We save the model checkpoint with the best validation performance. 
During inference, we reload this checkpoint and evaluate on a held-out test set (which is obtained from textual time series of unseen case reports).

For both tasks, we extract softmax probabilities rather than hard labels to allow for threshold calibration and ROC-based analyses. 
All training is conducted using PyTorch and HuggingFace Transformers.

%%%%%%%%%%%%%%%%%%%%%%%%%%%%%%%%%%%%%%%%%%%%%%%%%%%%%%%%%%%%%%%%%
%%%%%%%%%%%%%%%%%%%%%%%%%%%%%%%%%%%%%%%%%%%%%%%%%%%%%%%%%%%%%%%%
%\clearpage
\section{Survival Modeling with Language Model Embeddings}
\label{apd:encoder-survival}

To evaluate the prognostic value of textual information encoded in large language models (LLMs), we adopt a two-stage framework: (1) extraction of fixed-dimensional sequence embeddings from various pre-trained LLMs, and (2) downstream survival modeling using these embeddings as covariates.

In the first stage, we process each textual time series (which is converted to a textual context of \emph{"(time) clinical event [SEP] ... [SEP] (time) clinical event [SEP]"}) using a suite of LLMs to obtain dense vector representations that summarize the content of the input sequence. 
For models belonging to the encoder family, including \textit{bert-base-uncased}, \textit{roberta-base}, \textit{deberta-v3-small}, \textit{ModernBERT-base}, and \textit{ModernBERT-large}, we extract the final hidden state corresponding to the \texttt{[CLS]} token, which is conventionally used to represent the entire sequence in classification tasks. 
This token-specific embedding serves as a compact, sequence-level representation.

For decoder-based models that do not utilize a dedicated \texttt{[CLS]} token, including \textit{DeepSeek-R1-Distill-Llama-70B}, \textit{Llama-3.3-70B-Instruct}, \textit{DeepSeek-R1-Distill-Llama-8B}, and \textit{Llama-3.1-8B-Instruct}, we compute the mean-pooled embedding over the last hidden states of all non-padding tokens. 
This pooling strategy yields a fixed-length vector that captures the overall semantic content of the input while mitigating the impact of padding artifacts.

In the second stage, these embedding vectors are used as input covariates to three survival models, each designed to capture time-to-event dynamics in different ways:

\begin{itemize}
    \item \textbf{Random Survival Forest (RSF):} 
    A nonparametric ensemble method based on decision trees, capable of handling complex interactions and high-dimensional inputs. 
    The model is tuned over the number of trees (\texttt{n\_estimators}), the minimum number of samples required to split an internal node (\texttt{min\_samples\_split}), and the minimum number of samples at a leaf node (\texttt{min\_samples\_leaf}); see Table~\ref{tab:hyperparams} for search ranges.

    \item \textbf{DeepSurv:} 
    A neural-network generalization of the Cox proportional hazards model. 
    The architecture comprises fully connected layers with varying widths (\texttt{num\_nodes}) and dropout rates; see Table~\ref{tab:hyperparams} for search ranges. The number of training epochs is fixed.

    \item \textbf{DeepHit:} 
    A multi-task neural network that jointly models the discrete hazard and survival functions. 
    As with DeepSurv, we perform a grid search over hidden layer sizes and dropout rates while keeping the number of epochs fixed; see Table~\ref{tab:hyperparams}.
\end{itemize}

Hyperparameter tuning for each survival model is performed using a validation split from the training set. 
For each combination of language model and survival model, we conduct grid search over the respective hyperparameter space, selecting the configuration that maximizes the time-dependent concordance index \citep{antolini2005time} on the validation set. 
Final evaluation is conducted on held-out test sets (\textit{t2s2\_test}, \textit{sepsis-10}, \textit{sepsis-100}), and results are aggregated across survival metrics to assess model performance.

This framework allows for a systematic evaluation of how different pre-trained LLMs, and their associated embedding strategies, contribute to downstream survival prediction tasks.


\begin{table}[t]
\centering
\caption{Hyperparameter search ranges for survival models.}
\label{tab:hyperparams}
\begin{adjustbox}{max width=\columnwidth}
\begin{tabular}{ll}
\toprule
\textbf{Model} & \textbf{Hyperparameters} \\
\midrule
RSF & \texttt{n\_estimators} $\in$ \{20,50,100,200,500,1000\} \\
    & \texttt{min\_samples\_split} $\in$ \{2,10,20,50,100\} \\
    & \texttt{min\_samples\_leaf} $\in$ \{1,5,10,20\} \\
DeepSurv/DeepHit & \texttt{num\_nodes} $\in$ \{32,64,128,256,512,768,1024,1200,\\
& 1536,1800,2048,3072,4096\} \\
    & \texttt{dropout} $\in$ \{0.1,0.2,0.4,0.5\} \\
\bottomrule
\end{tabular}
\end{adjustbox}
\end{table}


\section{Computing Resources}
\label{apd:compute_resource}
All experiments were conducted on a high-performance computing cluster running AlmaLinux. 
For the largest language models (DSL33 and L33), training required three NVIDIA A100 GPUs 
(80~GB each; total 240~GB GPU memory), while inference was performed on two A100 GPUs. 
The encoder models were trained and evaluated on a single NVIDIA A100 GPU (80~GB). 
After getting the embedding vectors from LLMs, survival models were trained and evaluated on a workstation equipped with an NVIDIA RTX~4090 GPU. 
System memory usage was 32~GB RAM. 
Software dependencies, including specific versions of Python libraries and frameworks, are documented in \texttt{environment\_tts\_forecasting.yml} file of our \emph{anonymized} repository in the \emph{Supplementary Material}.
% : \url{https://shorturl.at/NHGBO}.


%%%%%%%%%%%%%%%%%%%%%%%%%%%%%%%%%%%%%%%%%%%%%%%%%%%%%%%%%%%%%%%%%
%%%%%%%%%%%%%%%%%%%%%%%%%%%%%%%%%%%%%%%%%%%%%%%%%%%%%%%%%%%%%%%%
%\clearpage
\onecolumn

{\refstepcounter{section} \label{TTS_prompt}
\Large\bfseries\thesection\quad LLM prompt to generate clinical textual time-series from PMOA case reports
}

% \begin{center}
\begin{minipage}{0.99\textwidth}
\begin{tcolorbox}[colback=gray!5, colframe=gray!100, title=Prompt, boxrule=2pt]

{You are a physician. 

Extract the clinical events and the related time stamp from the case report. The admission event has timestamp 0. If the event is not available, we treat the event, e.g. current main clinical diagnosis or treatment with timestamp 0. The events that happened before event with 0 timestamp have negative time, the ones after the event with 0 timestamp have positive time. The timestamp are in hours. The unit will be omitted when output the result. If there is no temporal information of the event, please use your knowledge and events with temporal expression before and after the events to provide an approximation. We want to predict the future events given the events happened in history. For example, here is the case report.}\\

\noindent
\textcolor{teal}{\texttt{An 18-year-old male was admitted to the hospital with a 3-day history of fever and rash. Four weeks ago, he was diagnosed with acne and received the treatment with minocycline, 100 mg daily, for 3 weeks. With increased WBC count, eosinophilia, and systemic involvement, this patient was diagnosed with DRESS syndrome. The fever and rash persisted through admission, and diffuse erythematous or maculopapular eruption with pruritus was present. One day later the patient was discharged.}}\\

\noindent
{Let's find the locations of event in the case report, it shows that four weeks ago of fever and rash, four weeks ago, he was diagnosed with acne and receive treatment. So the event of fever and rash happened four weeks ago, 672 hours, it is before admitted to the hospital, so the time stamp is -672. diffuse erythematous or maculopapular eruption with pruritus was documented on the admission exam, so the timestamp is 0 hours, since it happens right at admission. DRESS syndrome has no specific time, but it should happen soon after admission to the hospital, so we use our clinical judgment to give the diagnosis of DRESS syndrome the timestamp 0. then the output should look like:}\\
\vspace{2mm}
\noindent 

\texttt{18 years old | 0 \\
male | 0 \\
admitted to the hospital | 0 \\
fever | -72 \\
rash | -72 \\
acne | -672 \\
minocycline | -672 \\
increased WBC count | 0 \\
eosinophilia | 0 \\
systemic involvement | 0 \\
diffuse erythematous or maculopapular eruption | 0 \\
pruritis | 0 \\
DRESS syndrome | 0 \\
fever persisted | 0 \\
rash persisted | 0 \\
discharged | 24}\\

\noindent {Separate conjunctive phrases into its component events and assign them the same timestamp (for example, separation of ‘fever and rash’ into 2 events: ‘fever’ and ‘rash’). If the event has duration, assign the event time as the start of the time interval. Attempt to use the text span without modifications except ‘history of’ where applicable. Include all patient events, even if they appear in the discussion; do not omit any events; include termination/discontinuation events; include the pertinent negative findings, like ‘no shortness of breath’ and ‘denies chest pain’. Show the events and timestamps in rows, each row has two columns: one column for the event, the other column for the timestamp. The time is a numeric value in hour unit. The two columns are separated by a pipe $\mid$ as a bar-separated file. Skip the title of the table. Reply with the table only.}

\end{tcolorbox}
\end{minipage}
% \end{center}
\twocolumn

\onecolumn

{\refstepcounter{section}
\noindent\Large\bfseries\thesection\quad Prompting strategies for forecasting tasks
\label{apd:prompts}}

\subsection{F.1 \quad Zero-shot}

\begin{tcolorbox}[colback=gray!5, colframe=gray!100, title=Prompt, boxrule=2pt]

{You are an expert physician.}

{Reply to the prompt with structured predictions in a k-item, bar-separated row.

For example, if there k=3 events (A, B, C) and only B occurs in the time window, then:}
{0 $\mid$ 1 $\mid$  0
would be the correct response. \\ 

Reply with the answer only. Do NOT provide any other text. Do NOT explain.}

\end{tcolorbox}

\subsection{F.2 \quad Zero-shot (with degree of confidence)}

\begin{tcolorbox}[colback=gray!5, colframe=gray!100, title=Prompt, boxrule=2pt]
{You are an expert physician.}

{Reply to the prompt with structured predictions in a k-item, bar-separated row.
Assess whether the event happens in the time window using a scale from 1 to 5, where 5 means it definitely happened and 1 means it did not happen.

For example, if there k=3 events (A, B, C) and only B occurs in the time window, then:
$1 | 5 | 1$
would be the correct response. \\

Reply with the answer only. Do NOT provide any other text. Do NOT explain.}
\end{tcolorbox}

\subsection{F.3 \quad Few-shot}
\begin{tcolorbox}[colback=gray!5, colframe=gray!100, title=Prompt, boxrule=2pt]

{You are an expert physician.}

{Reply to the prompt with predictions in a k-item, bar-separated row to answer if the k events happen within a given forecast time window. \\

Example 1. If there k=3 events (A, B, C) and only B occurs in the requested time window, then:
0 $\mid$ 1 $\mid$ 0
would be the correct response.}\\

{Example 2. Suppose the time series is:}

\noindent \texttt{\textcolor{teal}
{admitted | 0 \\
history of diabetes | 0 \\
elevated troponin | 5 \\
--- \textbf{PREDICTION WINDOW} --- \\
cardiac catheterization | 6 \\
admission to icu | 24 \\ 
transfer to floor | 72\\ 
discharge | 144\\ 
office visit | 1440}
}

{Then, if the prediction is at 5 hours, the time preceding ``--- PREDICTION WINDOW ---'', the events admitted, history of diabetes, and elevated troponin have occurred.  Supposing k=4 and a forecast window of 24 hours, then the task it to determine which of cardiac catheterization, admission to icu, transfer to floor, and discharge happen. So the correct response for the example is:
1 $\mid$ 1 $\mid$ 0 $\mid$ 0
because cardiac catherization (6) and admission to icu (24) happen before 5 + 24 = 29.}\\

{Reply with the answer only. Do NOT provide any other text. Do NOT explain.}

\end{tcolorbox}

\twocolumn

\subsection{F.4 \quad Few-shot ordinal}

\begin{tcolorbox}[colback=gray!5, colframe=gray!100, title=Prompt, boxrule=2pt]

{You are an expert physician.}

{Reply to the prompt with predictions in a k-item, bar-separated row to answer if the k events happen within a given forecast time window, using a scale from 1 to 5, where 5 means it definitely happened and 1 means it did not happen.}\\

{Example 1. If there k=3 events (A, B, C) and only B occurs in the time window, then:
1 $\mid$ 5 $\mid$ 1
would be the correct response.}\\

{Example 2. Suppose the time series is:}\\
\noindent \texttt{\textcolor{teal}{
admitted | 0\\
history of diabetes | 0\\
elevated troponin | 5\\
--- \textbf{PREDICTION WINDOW} ---\\
cardiac catheterization | 6\\
admission to icu | 24\\
transfer to floor | 72\\
discharge | 144\\
office visit | 1440}}

{Then, if the prediction is at 5 hours, the time preceding ``--- PREDICTION WINDOW ---'', the events admitted, history of diabetes, and elevated troponin have occurred.  Supposing k=4 and a forecast window of 24 hours, then the task it to determine which of cardiac catheterization, admission to icu, transfer to floor and discharge happen. So the an excellent response for the example is:
5 $\mid$ 4 $\mid$ 1 $\mid$ 1
because cardiac catherization (6) and admission to icu (24) happen before 5 + 24 = 29.}\\

{Reply with the answer only. Do NOT provide any other text. Do NOT explain.}

\end{tcolorbox}

%%%%%%%%%%%%%%%%%%%%%%%%%%%%%%%%%%%%%%%%%%%%%%%%%%%%%%%%%%%%%%%%%
%%%%%%%%%%%%%%%%%%%%%%%%%%%%%%%%%%%%%%%%%%%%%%%%%%%%%%%%%%%%%%%%

\section{Open-source models for avoiding potential data leakage}
\label{apd:open_source}

We recognize that data leakage is a critical concern in evaluating model performance, particularly when using open-access text corpora. Below, we outline key factors that mitigate this concern in our study:

\textbf{1. Temporal annotations—not raw text—drive forecasting tasks.}
The core contribution of this work lies in the temporal structuring of clinical narratives. While large language models (LLMs) such as \textit{DeepSeek} and \textit{Llama-3.3} may have encountered PubMed Open Access (PMOA) case reports during pretraining, they were never exposed to the extracted timestamps or temporally ordered and future-masked (event, time) tuples used in our evaluations. These derived temporal annotations are essential to all forecasting tasks in our study—including F1-based next-event prediction, temporal ordering, and survival analysis—and constitute novel structured data not present in the original text.

\textbf{2. Temporal reordering fundamentally alters the learning problem.}
Even if LLMs were pretrained on PMOA narratives, they would have encountered events in their original, text-based order. In contrast, our framework chronologically reorders events based on extracted timestamps, creating a new temporal structure. This transformation introduces a distinct modeling challenge, as it breaks the sequential narrative flow seen during pretraining. As shown in Section 4.4.1, training on time-ordered versus text-ordered annotations leads to measurable differences in performance, underscoring the impact of this restructuring.

\textbf{3. External validation reveals limited effects of memorization.}
Our experimental design separates internal (T2S2 test set) and external validation (sepsis-10 and sepsis-100 cohorts), using entirely held-out case reports. The performance drop observed on external datasets (Tables 1–2) suggests that any potential memorization of PMOA case reports during pretraining yields minimal benefit and reinforces the robustness of our empirical findings.

\textbf{4. Broader context: standard practices and safeguards.}
Training data overlap is a general challenge in modern machine learning, especially for LLMs trained on web-scale corpora. Our work adheres to community standards while incorporating additional safeguards—such as explicit temporal annotation, event reordering, and external dataset validation—that reduce the risk of leakage and differentiate our tasks from standard text modeling.

To address concerns regarding potential pretraining overlap, we evaluated multiple models with documented training corpora. Both \textit{OLMO-32B-Instruct} (Allen Institute for AI) and \textit{RedPajama-INCITE-7B-Instruct} represent fully open-source models with transparent training data. Notably, RedPajama explicitly excludes PubMed from its training corpus, making it an ideal benchmark for assessing causal leakage. The comparable performance of these models to those potentially exposed to PMOA (see point 5 above) provides strong evidence that our temporal forecasting tasks are meaningfully distinct from standard language modeling on clinical narratives.

\textbf{5. Empirical evidence from PubMed-excluded models.}
Our results provide direct empirical evidence against meaningful causal leakage. \textit{RedPajama-INCITE-7B-Instruct}, which explicitly excludes PubMed from its training corpus, achieves comparable performance to other decoder models (0.352 F1 at 24h, 0.618 concordance) that potentially had access to PMOA during pretraining. This similarity in performance across models with and without PubMed exposure demonstrates that our T2S2 prediction task—with its temporal restructuring, event masking, and partial history inputs—is sufficiently distinct from raw case report text that prior exposure provides no meaningful advantage. Furthermore, \textit{OLMO-32B-Instruct} shows competitive performance (0.411 F1 at 24h), providing additional validation from a model with fully documented training data.

%%%%%%%%%%%%%%%%%%%%%%%%%%%%%%%%%%%%%%%%%%%%%%%%%%%%%%%%%%%%%%%%%
%%%%%%%%%%%%%%%%%%%%%%%%%%%%%%%%%%%%%%%%%%%%%%%%%%%%%%%%%%%%%%%%

\section{Broader Impact}
\label{apd:broader_impact}
This work advances equitable, scalable, and transparent medical AI by introducing a framework for forecasting clinical risk in sepsis from narrative documentation. Sepsis remains a leading cause of preventable mortality worldwide, with outcomes highly dependent on timely recognition and intervention. In many under‑resourced healthcare settings, structured EHR data and real‑time specialist input are unavailable, leaving unstructured narratives as the primary record of patient care. Our approach transforms such narratives into temporally structured representations, enabling systematic evaluation of forecasting methods for short‑ and long‑horizon event prediction, event ordering, and survival analysis. By releasing the dataset, extraction pipeline, and evaluation framework, we lower barriers for developing and benchmarking temporal forecasting models in sepsis, and promote reproducibility in LLM‑based clinical information extraction.

However, societal considerations are critical. LLM‑assisted extraction can introduce inaccuracies or hallucinations, and unvalidated forecasts risk amplifying bias or error. Although timelines are derived from non‑identifiable public case reports, their interpretation could still reinforce historical biases if not carefully contextualized. Forecasting results from case reports may also differ from those using routine clinical notes, limiting generalizability. Moreover, applying these forecasting models in high‑stakes clinical decision‑making without expert oversight poses safety risks. Responsible use requires transparent communication of dataset and model limitations, explicit uncertainty reporting, and safeguards to prevent inappropriate deployment in patient‑facing settings.
